\documentclass[11pt,a4paper]{article}
\usepackage[T1]{fontenc}
\usepackage{lmodern,microtype,amsmath,amssymb,amsthm,mathtools,booktabs}
\usepackage[margin=25mm]{geometry}
\usepackage[colorlinks=true,linkcolor=blue,citecolor=blue,urlcolor=blue]{hyperref}
\newtheorem{theorem}{Theorem}
\newtheorem{lemma}[theorem]{Lemma}
\newtheorem{corollary}[theorem]{Corollary}
\newcommand{\E}{\mathbb E}
\newcommand{\Prob}{\mathbb P}
\newcommand{\causal}{\mathrm{causal}}
\newcommand{\R}{\mathbb R}
\title{The sharp square-root bill of an exact causal source\\[4pt]
\large Full proof, audit bridge, and certified finite-horizon bounds}
\author{}
\date{3 October 2026}
\begin{document}
\maketitle
\begin{abstract}
Let $P,Q$ be finite probability laws of equal entropy $h$, with
surprisal standard deviations $a\ge b\ge0$. We prove
\[
C_n^{\causal}(P,Q)=nh+\sqrt{2/\pi}\,(a-b)\sqrt n+o(\sqrt n)
\]
for the full history-indexed exact causal model, including $b=0$.
The proof has two independent parts. A finite greedy response-tree
decomposition realizes the Bellman information envelope within
$\log_2 e$ bits at every horizon. An explicit two-Gaussian profile
approximates the centered Bellman recursion; Taylor's theorem and
sup-norm nonexpansiveness give a finite-horizon comparison with explicit
constants. Gaussian tail bounds then give convergence of the means.
No simultaneous coupling of policy diffusion limits, stochastic
differential equation convergence, or unproved recycling lemma is used.
For equal positive varentropies, a separate application of a published
nonuniform martingale Berry--Esseen bound yields an $O(\log(n+1))$ excess.
\end{abstract}

\section{The statement and the exact causal model}
Zero-probability symbols are removed. Entropies and surprisals are in
bits. Let
\[
H(P)=H(Q)=h,\qquad
a^2=\operatorname{Var}_P(-\log_2P(X)),\quad
b^2=\operatorname{Var}_Q(-\log_2Q(Y)),\quad a\ge b.
\]
A seed $S$ is sampled once, independently of external policy randomness.
At round $t$, its output is a deterministic function of $S$, the complete
past action-output history, and the currently requested action. Under
every adaptive policy, its conditional output law must be the requested
$P$ or $Q$. Let $C_n^{\causal}$ be the infimum entropy of finite seeds
with this property. Seeds may depend on $n$. Responses at distinct
histories are not required to be identical.

\begin{theorem}\label{thm:main}
For any such pair of finite laws,
\begin{equation}\label{eq:main}
\boxed{C_n^{\causal}(P,Q)
=nh+\sqrt{\frac2\pi}\,(a-b)\sqrt n+o(\sqrt n).}
\end{equation}
If $a=b=0$, then $C_n^{\causal}=nh$ exactly. If $a=b>0$, then also
\begin{equation}\label{eq:equalvar}
0\le C_n^{\causal}-nh=O(\log(n+1)).
\end{equation}
\end{theorem}

Write $\nu_P,\nu_Q$ for the laws of the centered one-letter surprisals
$D_A=-\log_2P_A(Y)-h$, so $\E D_A=0$. Define
\begin{equation}\label{eq:T}
(Tf)(x)=\min_{A\in\{P,Q\}}\E f(x-D_A),\qquad
G_0(x)=\mathbf1_{\{x\ge0\}},\qquad G_n=T^nG_0.
\end{equation}
Every $G_n$ is a finite-support CDF; let $Z_n\sim G_n$ and $e_n=\E Z_n$.
For a deterministic policy $\pi$, let $M_{\pi,n}$ be its transcript
surprisal minus $nh$. Backward induction gives the exact identity
\begin{equation}\label{eq:policy}
G_n(x)=\min_\pi\Prob(M_{\pi,n}\le x).
\end{equation}
Indeed, after choosing the first action and observing its output, each
continuation policy can be selected separately for the remaining
threshold. This is precisely \eqref{eq:T}. Allowing policy randomization
does not reduce the minimum: conditioning on its randomness expresses
its success probability as a mixture of deterministic-policy values.

The two ingredients of the proof are
\begin{equation}\label{eq:twoingredients}
0\le C_n^{\causal}-nh-e_n\le\log_2e,
\qquad
\frac{e_n}{\sqrt n}\longrightarrow\sqrt{2/\pi}\,(a-b).
\end{equation}
Both will be proved below.

\section{Global realization with a constant entropy gap}
We prove a finite-tree statement, also valid for history-dependent
kernels and any finite number of actions. A history is $v$, and its
child after action $A$ and output $y$ is $vAy$. Set
\[
p(\varnothing)=1,\qquad p(vAy)=p(v)K_v(y\mid A).
\]
These cylinder weights satisfy
$\sum_y p(vAy)=p(v)$ for each action separately. For a policy $\pi$,
let $\mathcal L_\pi$ denote its terminal histories. Their weights sum
to one and form its transcript law. Define $I_\pi=-\log_2p(T_\pi)$ and
\[
F_*(t)=\min_\pi\Prob(I_\pi\le t),\qquad Z\sim F_*,\quad L=\E Z.
\]

\begin{lemma}[Finite global realization]\label{lem:realization}
There is a finite exact causal seed $S_G$, supported on at most the
number of terminal histories in the full tree, such that
\begin{equation}\label{eq:gap}
L\le C\le H(S_G)\le L+\log_2e.
\end{equation}
\end{lemma}
\begin{proof}
A residual flow $r$ is nonnegative and satisfies
$\sum_y r(vAy)=r(v)$ for every action. A deterministic response
strategy $d$ chooses an output at each history-action pair. Its
compatibility indicator $\chi_d(v)$ equals one when all its outputs
on the path to $v$ agree with $v$. It is a flow of root mass one.

For any residual flow compute from the leaves upward
\[
B_r(v)=r(v)\quad(v\text{ terminal}),\qquad
B_r(v)=\min_A\max_y B_r(vAy)\quad(v\text{ nonterminal}).
\]
There are two elementary consequences.
First, the largest extractable strategy weight is $B_r(\varnothing)$:
for each possible action the strategy selects a child of maximal
bottleneck, and this is done recursively. Nonnegativity on terminal
histories implies nonnegativity at all histories by flow conservation.
Second, a policy choosing a minimizing action at every history has
\[
r(v)\le B_r(\varnothing)\qquad(v\in\mathcal L_{\pi_r}).
\]
Indeed, along every output of that policy, bottleneck values never
increase. This policy also proves that no larger strategy weight is
possible. If the root mass is positive, every action has a positive
child, and induction gives a positive root bottleneck.

Starting from $r=p$, repeatedly extract a maximizing strategy of weight
$w=B_r(\varnothing)$ and replace $r$ by $r-w\chi_d$. At least one
positive terminal residual becomes zero at each extraction; otherwise
the same strategy would permit a larger weight. Residuals never
increase, so the algorithm terminates after finitely many steps. Its
weights $w_j$ are nonincreasing, sum to one, and satisfy
\begin{equation}\label{eq:decomp}
p(v)=\sum_jw_j\chi_{d_j}(v).
\end{equation}
Choosing strategy $d_j$ with probability $w_j$ gives all the prescribed
history weights simultaneously. Ratios of child to parent weights
give the correct conditional laws at every positive-probability history.
Thus this is one exact causal seed for all policies.

Fix $0<\delta\le1$ and stop just after all weights larger than $\delta$
have been extracted. The remaining root mass is
$r(\varnothing)=\sum_{j:w_j\le\delta}w_j$. Its bottleneck is at most
$\delta$. The minimizing policy above therefore gives
\begin{align}
\sum_{j:w_j\le\delta}w_j
&\le\max_\pi\sum_{v\in\mathcal L_\pi}\min\{p(v),\delta\}\notag\\
&=\max_\pi\E\min\{1,\delta2^{I_\pi}\}
\le\E\min\{1,\delta2^Z\}.\label{eq:softtail}
\end{align}
The last inequality uses stochastic dominance: $F_*\le F_{I_\pi}$.
The zero residual case satisfies the same inequality trivially.
For $J=-\log_2w_{S_G}$, substitute $\delta=2^{-t}$ to obtain
\[
\Prob(J\ge t)\le\E\min\{1,2^{Z-t}\}
=\Prob(Z+E\ge t),
\]
where $E$ is independent exponential of rate $\ln2$. Integrating these
nonnegative-variable tails yields $H(S_G)\le L+1/\ln2$.

For the converse, under any deterministic policy each seed atom is
contained in its transcript atom. Its surprisal therefore dominates
the transcript surprisal under the induced coupling. Hence
$F_S\le F_{I_\pi}$ for all $\pi$, and thus $F_S\le F_*$. Integration
gives $H(S)\ge L$. This proves \eqref{eq:gap}.
\end{proof}

In the equal-entropy memoryless model, $F_*(t)=G_n(t-nh)$ by
\eqref{eq:policy}. Lemma~\ref{lem:realization} proves the first
inequality in \eqref{eq:twoingredients}. In particular, realizability
introduces no error proportional to the number of rounds or blocks.

\section{An explicit finite-horizon Bellman comparison}
Assume for now $a\ge b>0$. Let $\Phi$ be the standard normal CDF and set
\begin{equation}\label{eq:F}
F(x)=\begin{cases}
\displaystyle\frac{2b}{a+b}\Phi(x/b),&x\le0,\\[5pt]
\displaystyle1-\frac{2a}{a+b}\bigl(1-\Phi(x/a)\bigr),&x>0.
\end{cases}
\end{equation}
This is a CDF. Define
\[
K=\frac{\sqrt{2/\pi}}{a+b},\quad
m_3=\max_A\E|D_A|^3,\quad
\mathcal C=K\left(\frac{m_3}{3b^2}+a\right),\quad
u_t(x)=F(x/\sqrt t)\quad(t>0).
\]

\begin{lemma}[Explicit comparison]\label{lem:comparison}
For every $\tau\ge1$, $R\ge0$, and integer $n\ge0$, put
\[
\eta=\max\{F(-R/\sqrt\tau),\,1-F(R/\sqrt\tau)\}.
\]
Then, for every $x\in\R$,
\begin{equation}\label{eq:barrier}
\boxed{\begin{aligned}
F\!\left(\frac{x-R}{\sqrt{n+\tau}}\right)
-\eta-\frac{3\mathcal C}{\sqrt\tau}
&\le G_n(x),\\
G_n(x)&\le
F\!\left(\frac{x+R}{\sqrt{n+\tau}}\right)
+\eta+\frac{3\mathcal C}{\sqrt\tau}.
\end{aligned}}
\end{equation}
\end{lemma}
\begin{proof}
For $x$ of either sign let $c=b$ on the negative half-line and $c=a$
on the positive half-line. Direct differentiation gives
\[
F'(x)=Ke^{-x^2/(2c^2)},\qquad
F''(x)=-\frac{Kx}{c^2}e^{-x^2/(2c^2)}.
\]
The first two derivatives extend continuously across zero; in particular
$F''(0)=0$. The one-sided third derivatives satisfy
\[
|F'''(x)|=\frac K{c^2}
\left|\frac{x^2}{c^2}-1\right|e^{-x^2/(2c^2)}
\le\frac{2K}{b^2}.
\]
Consequently $F''$ is globally Lipschitz with that constant. Taylor's
theorem therefore applies even when an increment crosses the interface.

The function $u_t$ satisfies the following identity everywhere,
including $x=0$:
\begin{equation}\label{eq:pde}
\partial_tu_t(x)
=\frac12\min\{a^2\partial_{xx}u_t(x),b^2\partial_{xx}u_t(x)\}.
\end{equation}
On $x<0$ the second derivative is positive and the minimum selects
$b^2$; on $x>0$ it is negative and selects $a^2$. On each side both
expressions equal $-x\partial_xu_t(x)/(2t)$.

Because $\E D_A=0$, the spatial Taylor remainder gives
\[
\left|Tu_t-u_t-\partial_tu_t\right|
\le\frac{Km_3}{3b^2}\,t^{-3/2}
\qquad\text{uniformly in }x.
\]
For clarity, the minimum of the two quadratic Taylor terms is exactly
\eqref{eq:pde}; taking a minimum cannot increase the maximal absolute
remainder. Also, with $z=x/\sqrt t$,
\[
\partial_{tt}u_t(x)
=\frac K{4t^2}(3z-z^3/c^2)e^{-z^2/(2c^2)},
\qquad \|\partial_{tt}u_t\|_\infty\le\frac{2Ka}{t^2}.
\]
The last bound follows from
$\sup_{s\ge0}se^{-s^2/2}\le1$ and
$\sup_{s\ge0}s^3e^{-s^2/2}\le2$.
The time Taylor remainder is thus at most $Ka t^{-2}$. For $t\ge1$,
\begin{equation}\label{eq:local}
\|Tu_t-u_{t+1}\|_\infty\le\mathcal C t^{-3/2}.
\end{equation}

The operator $T$ is monotone, commutes with spatial translations and
addition of constants, and is nonexpansive in the sup norm. Telescoping
\eqref{eq:local} therefore gives
\begin{equation}\label{eq:telescoping}
\|T^nu_\tau-u_{n+\tau}\|_\infty
\le\mathcal C\sum_{j=0}^{n-1}(\tau+j)^{-3/2}
\le\frac{3\mathcal C}{\sqrt\tau}.
\end{equation}
Finally, the definition of $\eta$ implies the pointwise initial bounds
\[
u_\tau(x-R)-\eta\le\mathbf1_{\{x\ge0\}}
\le u_\tau(x+R)+\eta.
\]
Apply $T^n$, use its three invariances, and then
\eqref{eq:telescoping}. This proves \eqref{eq:barrier}.
\end{proof}

\begin{corollary}[CDF limit when $b>0$]\label{cor:cdf}
For each $x\in\R$, $G_n(x\sqrt n)\to F(x)$.
\end{corollary}
\begin{proof}
Fix $\tau,R$ and substitute $x\sqrt n$ in \eqref{eq:barrier}. Both
shifted arguments tend to $x$, so the limsup and liminf differ from
$F(x)$ by at most $\eta+3\mathcal C/\sqrt\tau$.
First choose $\tau$ large and then $R$ large. The error tends to zero.
For explicit choices at tolerance $0<\epsilon<1$, take
\[
\tau\ge\max\{1,(6\mathcal C/\epsilon)^2\},\qquad
R=a\sqrt{2\tau\ln(4/\epsilon)}.
\]
The Gaussian Chernoff bound gives $\eta\le\epsilon/2$, and the other
term is at most $\epsilon/2$. These parameters are fixed before sending
$n$ to infinity, so no uniform convergence over a moving policy grid
is being assumed.
\end{proof}

\section{Zero variance and convergence of the means}
\begin{lemma}[The degenerate endpoint]\label{lem:zero}
If $a>0$ and $b=0$, then
\[
G_n(x\sqrt n)\longrightarrow F_{a,0}(x)
=\begin{cases}0,&x\le0,\\2\Phi(x/a)-1,&x>0.\end{cases}
\]
\end{lemma}
\begin{proof}
This step concerns the centered increment control problem only.
For $\varepsilon>0$ add an independent Rademacher increment
$\varepsilon\xi_t$ at each time and for either action. Retain both the
original output and $\xi_t$ as observed labels. The new centered
increments have variances $a^2+\varepsilon^2$ and $\varepsilon^2$;
write $G_n^\varepsilon$ for their Bellman CDF.
Extra labels do not change the Bellman recursion: conditional on the
history the next increment law depends only on the action, and the
remaining terminal threshold depends only on the current sum.

The perturbed and unperturbed sums under the same action choices differ
by $\varepsilon\sum_{t=1}^n\xi_t$. For every $d>0$, Chebyshev gives
\[
\Prob\left(\left|\varepsilon\sum_{t=1}^n\xi_t\right|>d\sqrt n\right)
\le\varepsilon^2/d^2,
\]
uniformly over policies. Every original policy is a perturbed policy
that ignores the extra coins. Conversely, an original controller can
emulate a perturbed policy by drawing those coins as private randomness.
The optimal original value is no larger than that randomized policy's
value. These two observations give
\begin{equation}\label{eq:noise}
G_n^\varepsilon((x-d)\sqrt n)-\varepsilon^2/d^2
\le G_n(x\sqrt n)
\le G_n^\varepsilon((x+d)\sqrt n)+\varepsilon^2/d^2.
\end{equation}
Apply Corollary~\ref{cor:cdf} with
$a_\varepsilon=\sqrt{a^2+\varepsilon^2}$ and $b_\varepsilon=\varepsilon$.
Then let $\varepsilon\downarrow0$ with $d=\sqrt\varepsilon$.
The explicit CDFs \eqref{eq:F} converge to the continuous CDF $F_{a,0}$,
also at the moving arguments $x\pm\sqrt\varepsilon$. At $x=0$, both
limits are zero directly from the two formulas. Equation~\eqref{eq:noise}
proves the claim.
\end{proof}

\begin{lemma}[Mean limit]\label{lem:mean}
For $a\ge b\ge0$,
\begin{equation}\label{eq:mean}
\lim_{n\to\infty}\frac{e_n}{\sqrt n}
=\sqrt{\frac2\pi}\,(a-b).
\end{equation}
\end{lemma}
\begin{proof}
Let $B=\max_{A,y}|-\log_2P_A(y)-h|$. If $B=0$ the claim is immediate;
assume $B>0$. For any policy the centered
increments have conditional mean zero and absolute value at most $B$.
Convexity bounds their conditional exponential moment by
$\cosh(\lambda B)\le e^{\lambda^2B^2/2}$. Iteration and exponential
Markov therefore give, uniformly over all policies,
\[
\Prob(M_{\pi,n}>t\sqrt n),\quad
\Prob(M_{\pi,n}\le-t\sqrt n)
\le e^{-t^2/(2B^2)}\qquad(t>0).
\]
By \eqref{eq:policy}, both tails of $Z_n/\sqrt n$ have these same
bounds. Hence the signed tail formula
\[
\frac{e_n}{\sqrt n}
=\int_0^\infty[1-G_n(t\sqrt n)]\,dt
-\int_0^\infty G_n(-t\sqrt n)\,dt
\]
admits dominated convergence, using Corollary~\ref{cor:cdf} or
Lemma~\ref{lem:zero}. When $b>0$, its limit is
\begin{align*}
\int_0^\infty[1-F(t)]\,dt-\int_0^\infty F(-t)\,dt
&=\frac{2a^2}{(a+b)\sqrt{2\pi}}
-\frac{2b^2}{(a+b)\sqrt{2\pi}}\\
&=\sqrt{2/\pi}\,(a-b).
\end{align*}
Here $\int_0^\infty[1-\Phi(t/c)]\,dt=c/\sqrt{2\pi}$.
For $b=0<a$ the negative integral vanishes and the positive one is
$a\sqrt{2/\pi}$. If $a=b=0$, both increment laws are point masses
at zero and $G_n=G_0$, so the assertion is immediate.
\end{proof}

\begin{proof}[Proof of Theorem~\ref{thm:main}]
Lemma~\ref{lem:realization} and Lemma~\ref{lem:mean} give
\[
\frac{C_n^{\causal}-nh}{\sqrt n}
=\frac{e_n}{\sqrt n}+O(n^{-1/2})
\longrightarrow\sqrt{2/\pi}\,(a-b).
\]
If both varentropies vanish, both laws are uniform. Equal entropy means
their supports have the same size $k$. A single uniform seed on
$\{1,\ldots,k\}^n$, with the appropriate output relabeling for each
requested action, has entropy $n\log_2k=nh$. The all-$P$ transcript
gives the reverse inequality.
The stronger bound \eqref{eq:equalvar} is proved in the following section.
\end{proof}

\section{Equal positive varentropies: the logarithmic upper bound}
The main second-order argument above is self-contained. For the sharper
remainder when $a=b=\sigma>0$, we use the following specialization of
Fan--Grama--Liu, Theorem 2.1 \cite{FGL}.

\paragraph{External input (with its exact hypotheses).}
Let $(\xi_i,\mathcal F_i)$ be martingale differences such that, almost
surely,
\[
\sum_{i=1}^n\E(\xi_i^2\mid\mathcal F_{i-1})=1,\qquad
\left|\E(\xi_i^k\mid\mathcal F_{i-1})\right|
\le\tfrac12 k!\epsilon^{k-2}\E(\xi_i^2\mid\mathcal F_{i-1})
\quad(k\ge2),
\]
where $0<\epsilon\le1/2$. There is an absolute constant $C_{\rm FGL}$
such that, for all real $x$,
\begin{equation}\label{eq:fgl}
\left|\Prob\left(\sum_i\xi_i\le x\right)-\Phi(x)\right|
\le C_{\rm FGL}(1+x^2)\epsilon|\ln\epsilon|
\exp\left(-\frac{\widehat x_\epsilon^2}{2}\right),\qquad
\widehat x_\epsilon=\frac{2|x|}{1+\sqrt{1+2\epsilon|x|}}.
\end{equation}
This is that theorem with its quadratic-variation error parameter
$\delta=0$.

\begin{proof}[Proof of \eqref{eq:equalvar}]
Put $B=\max_{A,y}|-\log_2P_A(y)-h|$. For every deterministic policy use
\[
\xi_i=\frac{D_{A_i}(Y_i)}{\sigma\sqrt n},\qquad
\epsilon_n=\frac{B}{\sigma\sqrt n}.
\]
The conditional variance of every $\xi_i$ is exactly $1/n$, regardless
of which action the policy chooses. Also $|\xi_i|\le\epsilon_n$, so
\[
\left|\E(\xi_i^k\mid\mathcal F_{i-1})\right|
\le\epsilon_n^{k-2}\E(\xi_i^2\mid\mathcal F_{i-1})
\le\tfrac12 k!\epsilon_n^{k-2}\E(\xi_i^2\mid\mathcal F_{i-1}).
\]
Thus all the hypotheses hold once $n\ge4B^2/\sigma^2$, with identical
parameters for all policies. Taking the infimum of their CDFs preserves
the same two-sided error about $\Phi$. Consequently
\begin{equation}\label{eq:envfgl}
|G_n(\sigma\sqrt n\,x)-\Phi(x)|
\le C_{\rm FGL}(1+x^2)\epsilon_n|\ln\epsilon_n|
e^{-\widehat x_{\epsilon_n}^2/2}.
\end{equation}
This is a pointwise uniform bound, not merely a separate Wasserstein
estimate for each policy.

The right side has an integrable envelope independent of
$\epsilon_n\le1/2$. Indeed, for $|x|\ge1$,
\[
1+\sqrt{1+2\epsilon_n|x|}
\le1+\sqrt{1+|x|}\le3\sqrt{|x|},\qquad
\widehat x_{\epsilon_n}^2/2\ge2|x|/9.
\]
It follows that
\[
\int_{\R}(1+x^2)e^{-\widehat x_{\epsilon_n}^2/2}\,dx
\le\frac83+2\int_0^\infty(1+x^2)e^{-2x/9}\,dx<400.
\]
Since the standard normal has mean zero, the CDF formula for the
difference of means and \eqref{eq:envfgl} give
\begin{align*}
0\le e_n
&\le\sigma\sqrt n\int_{\R}|G_n(\sigma\sqrt n\,x)-\Phi(x)|\,dx\\
&\le400C_{\rm FGL}B\ln\left(\frac{\sigma\sqrt n}{B}\right).
\end{align*}
The nonnegativity follows also from stochastic dominance over any
single centered policy sum. Lemma~\ref{lem:realization} now gives,
for every integer $n\ge4B^2/\sigma^2$,
\[
0\le C_n^{\causal}-nh
\le400C_{\rm FGL}B\ln\left(\frac{\sigma\sqrt n}{B}\right)+\log_2e.
\]
Absorbing the finitely many smaller horizons into the constant proves
\eqref{eq:equalvar}.
\end{proof}

\section{The crash pair and the scope of the proof}
For $P=(1/2,1/4,1/4)$ and $Q=(q,q,1-2q)$ with
$q\in(1/3,1/2)$ defined by $H(Q)=3/2$, one obtains
\begin{equation}\label{eq:crash}
\boxed{C_n^{\causal}(P,Q)
=\tfrac32n+0.034747254051818\ldots\sqrt n+o(\sqrt n).}
\end{equation}
The exact coefficient is
$\sqrt{2/\pi}\bigl(1/2-\sqrt{V_Q}\bigr)$, where
\[
V_Q=2q(1-2q)\left(\log_2\frac q{1-2q}\right)^2.
\]
The exact-rational certificate in Section~\ref{sec:finite} encloses the
root, variance, and coefficient with explicit analytic remainders.
The rounded display in \eqref{eq:crash} is not the definition of the
coefficient.

\paragraph{What is repaired relative to the supplied draft.}
The proof uses the envelope over all adaptive transcript laws itself.
There is no need to identify a seed-dependent joint diffusion limit.
The discrete identity \eqref{eq:pde} has a minimum, and its sign is
verified on both half-lines. The finite-horizon inequality
\eqref{eq:barrier} replaces the unquantified diffusion/grid passage.
Finally, the single charge \eqref{eq:softtail} replaces the unproved
fiber-recycling estimate: it never accumulates a coder loss per block.

The draft's additional assertion $C_n^{\causal}=nh+O(\log n)$ when
$a=b>0$ does not follow from a vanishing square-root coefficient or a
fourth-moment bound. Section 5 proves it instead by an explicitly stated
external theorem, with all its hypotheses checked uniformly in the
policy. The full second-order formula \eqref{eq:main} and the stated
logarithmic upper bound are therefore both established. This does not
assert a sharp logarithmic coefficient or a logarithmic lower bound.

\paragraph{Verification and computational scope.}
The included exact verifier checks all residual flow identities,
reconstructed decoder marginals, and soft-tail inequalities in eighteen
finite instances, including history-dependent kernels and three
actions. It also includes a depth-three decoder certificate.
The separate Bellman-limit program provides floating-point
illustrations for $(a,b)=(2,1)$ and $(1,0)$; these numerical runs are
not used in the proof. The universal argument is the finite-tree proof
and the explicit inequalities above (and the stated external theorem for
the equal-variance refinement), not an extrapolation of tests.

The source and decoder may be exponential in $n$. There is no claim of
polynomial-time simulation, of one infinite-horizon seed, or of a
proof-assistant certificate. Literature priority has not been established.
Ordinary one-round coupling bounds and information-spectrum converses
are prior work \cite{C22,CKQGK23,SY23}; the response-tree realization
argument needed here is given in full in Lemma~\ref{lem:realization}.


\section{Reconciling the four audit requests}\label{sec:audit}
\paragraph{Lower bound.}
The threshold-policy-grid route in the earlier draft is replaced,
not completed by an implicit appeal to diffusion convergence.
Equation~\eqref{eq:policy} optimizes over every deterministic adaptive
policy at each finite horizon. Lemma~\ref{lem:comparison} compares that
exact finite-horizon value with the explicit profile \eqref{eq:F}.
Lemma~\ref{lem:mean} computes its mean and supplies a policy-uniform
integrable tail bound. Thus the mean limit does not rely on convergence
of a growing grid, or on interchanging a joint diffusion limit with
an infimum. Section~\ref{sec:finite} makes its finite-$n$ error explicit
for the crash pair.

\paragraph{Upper bound.}
The block/recycling route is likewise replaced. Lemma~\ref{lem:realization}
constructs a single exact response-tree seed and bounds its entire
entropy overhead by $\log_2 e$. No rounding of an arithmetic coordinate
and no antidiagonal lift are performed. In particular, charging a
constant per block of length $(\log n)^2$ would only give
$O(n/(\log n)^2)$, which is not $o(\sqrt n)$; the original route cannot
be certified by such an accounting argument. The global residual-flow
identity \eqref{eq:decomp} is the exactness certificate for the replacement.

\paragraph{Which profile gives one half?}
For the two constant policies alone, the centered normalized transcript
CDFs converge, when $a\ge b>0$, to $\Phi(x/a)$ and $\Phi(x/b)$.
Their lower envelope is
\[
F_{\rm fixed}(x)=
\begin{cases}\Phi(x/b),&x\le0,\\\Phi(x/a),&x>0.\end{cases}
\]
The same uniform Gaussian tail argument gives its mean
\[
\int_0^\infty\bigl(1-\Phi(t/a)\bigr)dt
-\int_0^\infty\Phi(-t/b)dt
=\frac{a-b}{\sqrt{2\pi}}.
\]
This is half of \eqref{eq:mean}. The full adaptive-policy envelope
is $G_n$, not the envelope of those two policies. Its limit is
\eqref{eq:F} and has twice this mean. Thus a statement that
\emph{every} policy-profile argument gives only half the coefficient
would be false in the model defined in Section~1.

The quoted assertion of the original ``Theorem 7'', namely that
marginal transcripts cannot yield more than $O(\sqrt n)$, does not
contradict a sharp positive multiple of $\sqrt n$. The present proof
identifies the coefficient within that order, and adds an achievability
argument using the compatible history tree. This comparison addresses
only the quoted assertion: the complete text and hypotheses of that
Theorem 7 have not been supplied with this audit, so no stronger
statement about its content is made.

\paragraph{What the numerical audit does.}
The following table encloses the \emph{actual optimum} using a proved
finite-$n$ inequality. It is not a direct numerical optimization of
$C_n$ at those horizons, nor a Bellman discretization passed off as an
exact seed cost. The executable checks the constants and endpoint
rounding using rational arithmetic. The analytic inequalities remain
an ordinary mathematical proof, not a proof-assistant certificate.

\section{An explicit error bound for the crash pair}\label{sec:finite}
Let $q$ be the unique root in $(1/3,1/2)$ of
\[
-2q\log_2q-(1-2q)\log_2(1-2q)=3/2.
\]
Uniqueness follows since the derivative is
$2\log_2((1-2q)/q)<0$ on this interval, and the endpoint entropies
are $\log_2 3$ and $1$. Put
\[
a=\tfrac12,\quad b=\sqrt{2q(1-2q)}\log_2\frac q{1-2q},
\qquad c=\sqrt{2/\pi}\,(a-b).
\]
These are exact definitions. The accompanying rational computation gives
\begin{align*}
0.410032428181985576080116712465
&<q<0.410032428181985576080116712466,\\
0.456450775263963373017231101124
&<b<0.456450775263963373017231101125,\\
0.034747254051817864080022137947
&<c<0.034747254051817864080022137948.
\end{align*}

\begin{theorem}[Quantitative crash-pair certificate]\label{thm:finite}
For every integer $n\ge16$,
\begin{equation}\label{eq:finite}
\boxed{\left|\frac{C_n^{\causal}-3n/2}{\sqrt n}-c\right|
\le \frac{17}{2}\frac{\sqrt{\ln n}}{n^{1/4}}
+\frac{3}{\sqrt{n\ln n}}+\frac{2}{\sqrt n}.}
\end{equation}
In particular the unnormalized remainder is
$O(n^{1/4}\sqrt{\ln n})$, with the displayed explicit bound.
\end{theorem}
\begin{proof}
The rational certificate verifies $9/20<b<1/2$ and $B<1$, where
$B$ is the increment bound in Lemma~\ref{lem:mean}. Since
$|D_A|^3\le B D_A^2$, it follows that $m_3\le1/4$. Also $\pi>2$,
so $K\le20/19$ and
\[
\mathcal C\le\frac{20}{19}
\left(\frac{1/4}{3(9/20)^2}+\frac12\right)
=\frac{4430}{4617}<1.
\]
For a given $n\ge16$ choose
\[
\tau=\sqrt n,\quad \eta_0=n^{-1/4},\quad
R=a\sqrt{2\tau\ln(2/\eta_0)},\quad
s=\sqrt{1+n^{-1/2}},\quad r=R/\sqrt n,\quad T=\sqrt{\ln n}.
\]
The Gaussian Chernoff bound and the prefactors at most two in
\eqref{eq:F} give $\eta\le\eta_0$. Thus, with
$d=\eta+3\mathcal C/\sqrt\tau\le4n^{-1/4}$,
Lemma~\ref{lem:comparison} reads
\[
F((x-r)/s)-d\le G_n(x\sqrt n)\le F((x+r)/s)+d.
\]
Let $Y=Z_n/\sqrt n$ and let $W$ have CDF $F$, so $\E W=c$.
Integrating the last comparison over $[-T,T]$ gives
\begin{equation}\label{eq:core}
\left|\int_{-T}^T\bigl[F(x/s)-G_n(x\sqrt n)\bigr]dx\right|
\le r+2Td.
\end{equation}
To see the constant $r$ explicitly, for any CDF $H$,
$\int_{\R}[H(x)-H(x-r)]dx=r$ when $r\ge0$. Apply this to $H(x)=F(x/s)$
for the upper signed bound and to $H(x+r)-H(x)$ for the lower one.
Each integrand is nonnegative, so restricting its integral to
$[-T,T]$ can only decrease it.

For $t>0$, Lemma~\ref{lem:mean} and \eqref{eq:F} give respectively
\[
\Prob(|Y|>t)\le2e^{-t^2/(2B^2)},\qquad
\Prob(|sW|>t)\le2e^{-t^2/(2a^2s^2)}.
\]
The latter follows by bounding both Gaussian scales by $a$; their
prefactors sum to two. For any $v>0$,
$\int_T^\infty e^{-t^2/(2v)}dt\le(v/T)e^{-T^2/(2v)}$.
Since $B^2\le1$, $a^2s^2\le1$, and
$2B^2+2a^2s^2\le3$, the total contribution outside $[-T,T]$
to the absolute signed difference of the means is at most
\[
\frac3T e^{-T^2/2}=\frac{3}{\sqrt{n\ln n}}.
\]
Consequently \eqref{eq:core} bounds $|\E Y-sc|$ by
$r+2Td+3/\sqrt{n\ln n}$. Next,
$0\le(s-1)c\le1/(2\sqrt n)$, using $c<1$ and
$\sqrt{1+u}-1\le u/2$. For $n\ge16$,
$\ln2\le\tfrac14\ln n$, whence
\[
r=\tfrac12 n^{-1/4}\sqrt{2\ln(2n^{1/4})}
\le\tfrac12 n^{-1/4}\sqrt{\ln n},\qquad
2Td\le8n^{-1/4}\sqrt{\ln n}.
\]
Finally, Lemma~\ref{lem:realization} changes $\E Y$ to
$(C_n^{\causal}-3n/2)/\sqrt n$ by at most
$\log_2e/\sqrt n<3/(2\sqrt n)$; the last inequality follows from
$\ln2>2/3$. Adding these bounds proves \eqref{eq:finite}.
\end{proof}

\begin{center}
\begin{tabular}{@{}rcc@{}}
\toprule
$n$ & Certified lower endpoint & Certified upper endpoint\\
\midrule
$10^{8}$ & 0.000000000000 & 0.399831377112 \\
$10^{12}$ & 0.000000000000 & 0.079430259815 \\
$10^{16}$ & 0.029587976869 & 0.039906531235 \\
$10^{20}$ & 0.034170431871 & 0.035324076233 \\
$10^{24}$ & 0.034684066372 & 0.034810441732 \\
$10^{32}$ & 0.034746524423 & 0.034747983681 \\
$10^{48}$ & 0.034747253962 & 0.034747254142 \\
\bottomrule
\end{tabular}
\end{center}
Every row encloses $(C_n^{\causal}-3n/2)/\sqrt n$. Endpoints are rounded
outward to twelve decimal places. Nonnegativity is also used for the
first two lower endpoints. Already the row $n=10^{16}$ excludes
$c/2=0.0173736270259089\ldots$; later rows narrow around $c$.

\paragraph{Reproduction and arithmetic certificate.}
Run \texttt{python3 certify\_crash\_pair.py}. Only the Python standard
library is required. The output JSON records the coefficient intervals,
all table rows, and the status \texttt{PASS\_EXACT\_RATIONAL\_ARITHMETIC}.
Logarithms are enclosed with 64 terms of the positive series
\[
\ln x=2\sum_{j=0}^{63}\frac{z^{2j+1}}{2j+1}+R,
\qquad z=\frac{x-1}{x+1},\quad
0\le R\le\frac{2z^{129}}{129(1-z^2)}\quad(1\le x\le2),
\]
after powers-of-two range reduction. Entropy signs are checked exactly
at both endpoints at every root-bracketing stage. Square roots use
integer-square-root enclosures; $\pi$ uses Machin's identity with
alternating-series remainders. Every decimal enclosure is rounded
outward. No floating-point value is accepted as an arithmetic premise.
The script certifies this arithmetic, not the universal analytic theorem
by formal verification. The standalone proof above supplies that theorem.

\begin{thebibliography}{9}
\bibitem{C22} S. Compton.
\emph{A Tighter Approximation Guarantee for Greedy Minimum Entropy
Coupling}. 2022. \href{https://arxiv.org/abs/2203.05108}{arXiv:2203.05108}.
\bibitem{CKQGK23} S. Compton, D. Katz, B. Qi, K. Greenewald, and M. Kocaoglu.
\emph{Minimum-Entropy Coupling Approximation Guarantees Beyond the
Majorization Barrier}. AISTATS, PMLR 206, 10445--10469, 2023.
\href{https://proceedings.mlr.press/v206/compton23a.html}{Proceedings}.
\bibitem{SY23} Y. Y. Shkel and A. K. Yadav.
\emph{Information Spectrum Converse for Minimum Entropy Couplings and
Functional Representations}. ISIT, 2023.
\href{https://arxiv.org/abs/2305.05745}{arXiv:2305.05745}.
\bibitem{FGL} X. Fan, I. Grama, and Q. Liu.
\emph{Nonuniform Berry--Esseen bounds for martingales with applications
to statistical estimation}. Theorem 2.1, conditions (A1)--(A2).
\href{https://arxiv.org/abs/1608.05217}{arXiv:1608.05217}.
\end{thebibliography}
\end{document}
